\paragraph{A common unipotent group and arithmetic periodization.}
Let \(E/F\) be a CM extension and \((V,h)\) a Hermitian space,
with \(h\) conjugate-linear in its first argument. Fix hyperbolic pairs
\(e_i,f_i\), \(1\le i\le r\), put \(X=\langle e_i\rangle\),
and let the vectors \(z_u\) lie in
\(W=\langle e_i,f_i\rangle^\perp\). The apparently different groups
\(N_u\) in the period construction are all the same group:
\[
 N=\{n\in U(V): (n-1)e_i\in\langle e_j:j<i\rangle,
                  \ (n-1)X^\perp\subset X\}.
\]
Indeed, their additional condition \((n-1)z_u\in X\) follows from
\(z_u\in X^\perp\). This argument also holds over the original finite
adelic coefficient algebra. The characters can still vary with \(u\):
for \(r\ge1\), their ratio is
\[
 \chi_u(n)\overline{\chi_v(n)}
 =\Psi_E\!\left(m^{-1}h(f_r,n(z_u-z_v))\right),
\]
where \(m>0\) is the level denominator. Distinct characters are
orthogonal on the compact quotient \(N(F)\backslash N(\mathbb A_F)\)
with probability measure. This gives Fourier orthogonality there;
repeated characters and the other rational orbits remain to be treated.
For \(r=0\), the group and its character are trivial.

Keep the original normalizing constants
\[
 c_u=\frac{v_Kv_{N,u}v_{H,u}}{d_u}>0,\qquad d_u=\dim\pi_u,
\]
where \(v_K\) is the finite adelic level volume, \(v_{N,u}\) the
archimedean unipotent lattice volume, and \(v_{H,u}\) the compact
archimedean stabilizer volume. A common Haar measure on \(N_\infty\)
is compatible with the normalized maps: if
\(\mu_{N,u}=t_u\mu_N\), the new quotient measure is \(t_u\nu_u\),
the unit vector becomes \(t_u^{-1/2}W_u\), and
\(P'_u=t_u^{-1/2}P_u\), \(c'_u=c_u/t_u\). Thus
\(P'_u/\sqrt{c'_u}=P_u/\sqrt{c_u}\). The arithmetic maps, their
integral lattices and metrics remain the original ones.

\begin{proposition}[Periodization on the actual seed spaces]
\label{shared-common-periodization}
Fix finitely many words. Write \(G=U(V)(\mathbb A_F)\),
\(\Gamma=U(V)(F)\), and let \(\Delta_u\subset\Gamma\) be their
rational stabilizers. Use compatible quotient measures on
\(Y=\Gamma\backslash G\) and \(X_u=\Delta_u\backslash G\).
Let \(C_u\) be the original finite-dimensional Hilbert column spaces
and \(\Phi_u:C_u\to L^2(X_u)\) their linear maps to scalar,
left-invariant seeds. Write
\(P_ua(y)=\sum_{\delta\in\Delta_u\backslash\Gamma}
\Phi_u(a)(\delta s(y))\).
Assume, for each basis vector \(a\), that
\(\sum_{\delta\in\Delta_u\backslash\Gamma}
 |\Phi_u(a)(\delta s(y))|\) belongs to \(L^2(Y)\), where \(s\)
is a measurable section. Then
\[
 \bigoplus_u L^2(X_u)\simeq
 L^2\!\left(Y,\bigoplus_u\ell^2(\Delta_u\backslash\Gamma)\right),
 \qquad f\longmapsto\bigl(f_u(\delta s(y))\bigr)_{u,\delta}
\]
is unitary, and the normalized synthesis is
\[
 Sa=\sum_uc_u^{-1/2}P_ua_u=\mathcal F Aa,\qquad
 (Aa)_{u,\delta}(y)=\frac{\Phi_u(a_u)(\delta s(y))}{\sqrt{c_u}},
 \quad \mathcal Fv=\sum_{u,\delta}v_{u,\delta}.
\]
Here \(\mathcal F\) is used on the domain where the sum of absolute
values belongs to \(L^2(Y)\). Under absolute unfolding, its full norm is
\begin{equation}\label{shared-periodization-orbits}
 \|Sa\|^2=\sum_{u,v}\frac1{\sqrt{c_uc_v}}
 \sum_{\gamma\in\Delta_u\backslash\Gamma/\Delta_v}
 \int_{(\Delta_u\cap\gamma\Delta_v\gamma^{-1})\backslash G}
 \Phi_u(a_u)(g)\overline{\Phi_v(a_v)(\gamma^{-1}g)}\,dg.
\end{equation}
\end{proposition}
\begin{proof}
The sets \(\delta s(Y)\) partition a fundamental region for
\(\Delta_u\). Integration over this partition proves the isometry
and constructs its inverse. Summing the seed values gives the stated
synthesis. Expanding its squared norm and partitioning the rational
pairs into double cosets gives \eqref{shared-periodization-orbits};
absolute unfolding justifies every rearrangement. For seeds covariant
under a unitary representation, the same construction uses its
associated bundle and unitary cocycle. The scalar seeds in the statement
already include the compact coefficient contraction and are invariant.
\end{proof}

The unitary identification does not make \(A\) an isometry without
the additional identity
\(\int_{X_u}\Phi_u(a)\overline{\Phi_u(b)}=c_u\langle a,b\rangle\)
for the actual column metric. Nor is unrestricted fibre summation
bounded when \([\Gamma:\Delta_u]=\infty\). This index is infinite
for \(r\ge1\): the rational unitary maps
\(e_1\mapsto e_1+t f_1\), with \(\bar t=-t\), fixing \(f_1\)
and the orthogonal complement, give distinct cosets because
\(\Delta_u\) fixes \(e_1\). The estimate needed concerns the
finite seed spaces with their prescribed restrictions.

\begin{corollary}[A calibrated summation criterion]
\label{shared-periodization-calibration}
Let \(J_u:C'_u\hookrightarrow C_u\) be the actual isometric source
inclusions and \(\Pi\) the orthogonal projection onto the final target
\(H_0\subset L^2(Y)\), of dimension \(n>0\). Put
\[
 Q_0=\sum_uc_u^{-1}\Pi P_uJ_uJ_u^*P_u^*\Pi,
 \qquad M_0=\operatorname{tr}Q_0>0.
\]
Suppose positive measurable weights satisfy, for all \(a_u\in C'_u\),
\begin{equation}\label{shared-periodization-weights}
 \sum_{u,\delta}w_{u,\delta}(y)\le L,\qquad
 \int_Y\sum_{u,\delta}
 \frac{|\Phi_u(J_ua_u)(\delta s(y))|^2}{c_uw_{u,\delta}(y)}\,dy
 \le B\sum_u\|a_u\|^2,
\end{equation}
where \(L,B>0\), and the first inequality holds almost everywhere.
If \(LB\le(1+\varepsilon)M_0/n\), with \(\varepsilon\ge0\), then
\[
 Q_0\le LB I_{H_0},\qquad
 \frac{\operatorname{tr}(Q_0-(M_0/n)I_{H_0})_+}{M_0}
 \le\frac{\varepsilon}{1+\varepsilon}.
\]
\end{corollary}
\begin{proof}
Weighted Cauchy--Schwarz gives
\(\|\mathcal F AJ a\|^2\le LB\|a\|^2\); applying \(\Pi\)
does not increase the norm. Thus every eigenvalue \(\lambda\) of
\(Q_0\) lies in \([0,LB]\), and
\((\lambda-M_0/n)_+\le(1-M_0/(nLB))\lambda\).
Summing proves the trace bound. The mass assumption implies
\(LB\ge M_0/n\).
\end{proof}

Constructing weights with this calibration, or proving the corresponding
summed norm bound directly, remains an arithmetic estimate. All the
orbits in \eqref{shared-periodization-orbits}, the fixed normalizing
volumes, and the arithmetic height hypotheses must still be respected.
